% =========================================================
% Proof: Classification of Euclidean Rigid Motions
% Source: notes/isometries/notes-isometries.tex
% Difficulty: Hard — requires linear algebra and careful case analysis
% =========================================================

% ⚠️  DIFFICULTY WARNING
% This proof requires:
%   - Knowing that a rigid motion fixing the origin must be linear
%   - Showing linear isometries of R^n are exactly orthogonal maps
%   - Decomposing any rigid motion as (orthogonal) ∘ (translation)
% Recommended: attempt only after you are comfortable with all other
% isometry proofs and have reviewed orthogonal matrices in linear algebra.

\subsection*{Classification of Euclidean Rigid Motions}
\label{prf:classification-rigid-motions}

\begin{remark*}[Return]
$\leftarrow$ \hyperref[thm:classification-rigid-motions]{Back to Theorem in Notes}
\end{remark*}

\begin{proof}
\Claim Every rigid motion $f : \mathbb{R}^n \to \mathbb{R}^n$ can be written
uniquely as $f = T_a \circ R$ where $R$ is orthogonal and $a = f(0)$.

\Given A bijective isometry $f : \mathbb{R}^n \to \mathbb{R}^n$.

\Goal To show $f$ decomposes as translation composed with orthogonal map.

\Strategy Let $a = f(0)$ and define $g = T_{-a} \circ f$, so $g(0) = 0$.
Show $g$ is an isometry fixing the origin, hence linear, hence orthogonal.
The key lemma is: an isometry of $\mathbb{R}^n$ fixing $0$ must be linear.

% -------------------------------------------------------
% YOUR PROOF HERE
% -------------------------------------------------------

\end{proof}

\begin{remark*}[Proof shape]
Reduce to the origin-fixing case, then use the polarization identity
to show linearity: $d(f(x), f(y)) = d(x,y)$ and $f(0) = 0$ together
force $f(x+y) = f(x) + f(y)$ and $f(\lambda x) = \lambda f(x)$.
\end{remark*}

\begin{remark*}[Dependencies]
\begin{itemize}
  \item Proposition~\ref{prop:translations-are-isometries}: translations.
  \item Proposition~\ref{prop:rotations-are-isometries}: orthogonal maps.
  \item Linear algebra: inner product via polarization,
        $\langle x,y \rangle = \tfrac{1}{2}(\|x\|^2 + \|y\|^2 - \|x-y\|^2)$.
  \item A linear isometry of $\mathbb{R}^n$ is represented by an orthogonal matrix.
\end{itemize}
\end{remark*}
